%%% Section: Work Experience

\sectionTitle{Профессиональный опыт}{\faSuitcase}
\begin{experiences}
  \experience[img/logo/mws]
    {наст. время} {Kotlin-разработчик}{\link{https://mws.ru}{МТС Веб Сервисы}}{Санкт-Петербург, Россия}
    {Ноябрь 2025}  {
                     \begin{itemize}
                     \item Команда Kotlin API and Gateway, IT-кластер <<Облачные технологии>> (Development Platform);
                     \item Развивал внутреннюю Kotlin-платформу и API-шлюз для команд облачных сервисов;
                     \item Разрабатывал инструмент кодогенерации, который по OpenAPI-спецификациям сервисов строит модели, клиентов и код межсервисного взаимодействия;
                     \item Спроектировал и написал сервис автоматизации релизов и управления проектами: планирование сборок, релизные ветки и кросс-репозиторная выкатка; интеграции с GitLab, Jira и Telegram;
                     \item Реализовал маскирование конфиденциальных данных в access-логах по типам запросов и ответов с минимальным оверхедом.
                     \end{itemize}
                  }
                  {Kotlin, Spring Boot, Spring WebFlux, R2DBC, Gradle, OpenAPI, Shell, GitLab CI/CD, Docker, Kubernetes, Helm}
  \emptySeparator
  \experience[img/logo/itmo-ru]
    {наст. время} {Преподаватель практики}{\link{https://itmo.ru}{Университет ИТМО}}{Санкт-Петербург, Россия}
    {Сентябрь 2025}  {
                     \begin{itemize}
                     \item Вёл практику у 54 студентов ПИиКТ по дисциплинам <<Распределённые системы хранения данных>>, <<Тестирование программного обеспечения>> и <<Веб-программирование>>;
                     \item Проверял лабораторные работы, разбирал ошибки в коде; принимал защиты и оценивал рубежные работы;
                     \item Составлял дополнительные задания, объяснял сложные темы и консультировал студентов.
                     \end{itemize}
                  }
                  {}
  \emptySeparator
  \experience[img/logo/followy]
    {Декабрь 2021} {Frontend-разработчик}{\link{https://vk.com/followy}{Followy}}{Ростов-на-Дону, Россия}
    {Август 2021}  {
                     \begin{itemize}
                     \item Разрабатывал веб-клиент Followy: адаптивный и доступный (a11y) пользовательский интерфейс и модульные тесты (Jest, Enzyme) для обеспечения надёжности;
                     \item Реализовал парсинг и компиляцию внутреннего языка разметки в приложении;
                     \item Профилировал производительность приложения (Lighthouse), устраняя узкие места.
                     \end{itemize}
                  }
                  {JavaScript, TypeScript, React.js, Redux, Material UI, Jest, Docker, PWA, Sentry}
\end{experiences}
